% !TeX program = xelatex
\documentclass[UTF8,fontset=fandol,zihao=-4,a4paper]{ctexart}
\usepackage[top=25mm,bottom=26mm,left=26mm,right=26mm,headheight=15pt]{geometry}
\usepackage{mathtools,amssymb,amsthm,bm}
\usepackage{booktabs,longtable,array,enumitem,microtype}
\usepackage{fancyhdr,xcolor,hyperref}
\setmainfont{TeX Gyre Termes}
\hypersetup{unicode=true,colorlinks=true,linkcolor=black,citecolor=black,
  urlcolor=blue!45!black,pdftitle={三元平均的有限步可达性：算术判据与重复值不变量},
  pdfauthor={研究整理稿},pdfsubject={有理配置的有限次三元等权平均}}
\setlength{\parskip}{0.28em}
\setlength{\emergencystretch}{2em}
\setlist{itemsep=0.25em,topsep=0.4em,leftmargin=2em}
\numberwithin{equation}{section}
\newtheorem{theorem}{定理}[section]
\newtheorem{lemma}[theorem]{引理}
\newtheorem{proposition}[theorem]{命题}
\newtheorem{corollary}[theorem]{推论}
\theoremstyle{definition}
\newtheorem{definition}[theorem]{定义}
\newtheorem{remark}[theorem]{注}
\newtheorem{example}[theorem]{例}
\renewcommand{\proofname}{证明}
\newcommand{\Z}{\mathbb Z}
\newcommand{\Q}{\mathbb Q}
\newcommand{\R}{\mathbb R}
\newcommand{\F}{\mathbb F}
\newcommand{\PP}{\mathcal P}
\newcommand{\HH}{\mathcal H}
\newcommand{\II}{\mathcal I}
\newcommand{\B}{\mathcal B}
\newcommand{\one}{\boldsymbol 1}
\DeclareMathOperator{\prim}{prim}
\DeclareMathOperator{\diag}{diag}
\DeclareMathOperator{\lcm}{lcm}
\DeclareMathOperator{\sgn}{sgn}
\newcommand{\mat}[4]{\begin{pmatrix}#1&#2\\#3&#4\end{pmatrix}}
\pagestyle{fancy}
\fancyhf{}
\fancyhead[L]{\small 三元平均的有限步可达性}
\fancyhead[R]{\small 研究稿}
\fancyfoot[C]{\thepage}
\renewcommand{\headrulewidth}{0.3pt}
\title{\heiti 三元平均的有限步可达性\\[0.3em]
  \large 算术判据与重复值不变量}
\author{}
\date{2026 年 9 月 11 日\quad 研究整理稿（补充修订）}

\begin{document}
\maketitle
\thispagestyle{plain}
\begin{abstract}
给定有限个有理数，一次三元平均选取三个不同位置，并将其数值同时替换为三者的算术平均。
本文研究在位置总数固定、不增加副本的条件下，何时可以有限步达到全体相等。
对非恒定输入，减去总体均值、清除分母并取本原整数代表
$X=(x_1,\ldots,x_n)$，定义 $G(X)=\gcd_{i,j}(x_i-x_j)$。
我们证明：当 $n\ge7$ 时，可达性等价于 $G(X)$ 是 $3$ 的非负整数次幂。
证明的主体直接处理所有 $n\ge11$，不区分素数与合数，也不调用较小维数的结论。
其方法是预留一位三进精度，建立两个不同三重值的不变量，并利用危险支持集、仿射同余见证和整数平方能量证明有限终止。
对整数零和合法输入，$n\ge11$ 时全部中间值可限制在 $\frac13\Z$。
七、八、十元分别由独立的低维构造处理，九元使用三进张量网络；其中七、十元所需的精确有限代数证书在附录中列出。
当 $n\ge7$ 时，判定算法关于有理输入的总位长为多项式，但本文不证明这些维数具有统一的位长多项式构造长度。
小维数方面，四元具有直接充要条件；五元具有完整确定性递归，步数为 $O((\log(2+H))^2)$，因此判定及成功路径构造均为位长多项式；六元具有 $O(H^3)$ 状态规模的完备整数图判定。
五、六元的反例说明统一算术判据的维数条件不可删除。
\end{abstract}
\noindent\textbf{关键词：}有限步平均；混合图；有理数可达性；同余不变量；重复值；计算机辅助证明。

\tableofcontents
\newpage

\section{问题、主要结论与证明结构}

设 $n\ge3$。对 $x=(x_1,\ldots,x_n)\in\Q^n$ 及三个不同的下标
$i,j,k$，允许的操作为
\begin{equation}\label{eq:atom}
 (x_i,x_j,x_k)\longmapsto(m,m,m),\qquad
 m=\frac{x_i+x_j+x_k}{3},
\end{equation}
其他位置不变。操作次数必须有限；三个位置可在每一步根据当前状态重新选择。
各步保持总和，故终态若为常值，只能等于初始均值
$\mu=n^{-1}\sum_i x_i$。
我们始终使用原来的 $n$ 个位置；相等数值的多份出现仍代表不同位置。

对非恒定有理输入，先作中心化，再选正有理数 $c$ 使
\begin{equation}\label{eq:normalization}
 X=c(x-\mu\one)\in\Z^n,\qquad
 \sum_i X_i=0,\qquad \gcd_i X_i=1.
\end{equation}
这样的 $X$ 在整体符号之外唯一。定义
\begin{equation}\label{eq:G}
 G(X)=\gcd_{1\le i<j\le n}|X_i-X_j|>0.
\end{equation}

\begin{theorem}[完整算术判据]\label{thm:main}
设 $n\ge7$，输入为 $n$ 个有理数。若输入已全相等，则无需操作。
否则取 \eqref{eq:normalization} 的本原中心化代表 $X$。
输入可由有限次 \eqref{eq:atom} 变为常值，当且仅当
\[
 G(X)=3^a\qquad\text{对某个整数 }a\ge0.
\]
\end{theorem}

下面的统一结论是本文主体。
\begin{theorem}[全维数的直接构造与精度界]\label{thm:uniform}
设 $n\ge11$，$X\in\Z^n$ 的坐标和为零。若对每个非三素因子
$q\mid n$，坐标模 $q$ 不全同余，则 $X$ 可在原 $n$ 个位置上有限步归零，
且全部中间值属于 $\frac13\Z$。

更一般地，只要非零 $X$ 的本原代表满足 $G=3^a$，同一精度结论成立。
\end{theorem}

第 \ref{sec:arithmetic} 节给出必要条件及统一的算术语言。
第 \ref{sec:safety}--\ref{sec:termination} 节直接证明定理 \ref{thm:uniform}。
这些章节不依赖小维数定理、素数归约、合数乘法闭包或抽象矩阵的实现定理。
第 \ref{sec:small} 节处理 $7\le n\le10$，从而完成定理 \ref{thm:main}。
附录 \ref{app:seven}、\ref{app:ten} 记录其中两个精确有限证书及其验证规则。
第 \ref{sec:boundary} 节补齐四、五、六元的结果：四元的显式条件，
五元的递归充要判据及平方对数步数界，以及六元的完整伪多项式判定。

二元平均的对应问题由 Coviello Gonzalez 与 Chrobak \cite{mixability} 研究。
其充分性证明在预留一位二进精度后维持两个不同的重复值。
本文借鉴这一证明结构；三元情形需要同时处理三个模 $3$ 余类、
每次生成三份新值的重数变化，以及模 $2$ 下不能随意除以 $2$ 的问题。
定理 \ref{thm:uniform} 是这里给出的三元论证，不是对 \cite{mixability} 中定理的直接引用。

\section{算术不变量与归一化}\label{sec:arithmetic}

\begin{lemma}[仿射相容性]\label{lem:affine-equivariance}
固定任意一个位置操作序列 $W$。对 $c\ne0$、$b\in\Q$，有
\[
 W(cx+b\one)=cW(x)+b\one.
\]
因此中心化、整体非零缩放、坐标重新编号均不改变归零可达性。
\end{lemma}
\begin{proof}
单次平均满足该恒等式，依操作次数归纳即可。坐标重新编号只改变之后选择的位置编号，
不构成额外的数值操作。
\end{proof}

\begin{lemma}\label{lem:G-divides-n}
若 $X\in\Z^n$ 非零、本原且零和，则 $G(X)\mid n$。
此外，$G(X)=3^a$ 等价于：对每个素数 $q\mid n$、$q\ne3$，
$X$ 模 $q$ 不全同余。
\end{lemma}
\begin{proof}
记 $G=G(X)$。由 $G\mid X_i-X_1$ 得 $G\mid nX_1$。
又 $\gcd(G,X_1)=1$，否则某素数同时整除所有坐标，与本原性矛盾。
故 $G\mid n$。最后，对任意素数 $q$，坐标模 $q$ 全同余恰当且仅当 $q\mid G$。
\end{proof}

\begin{proposition}[必要性]\label{prop:necessary}
任意维数中，非零本原零和整数状态若能归零，则 $G(X)$ 为三幂。
\end{proposition}
\begin{proof}
若 $q\ne3$ 为 $G(X)$ 的素因子，则所有坐标模 $q$ 同余于某个 $r\ne0$；
这里 $r\ne0$ 来自本原性。记 $R=\Z[1/3]$。
所有中间值都在 $R$ 中，且因 $3$ 在 $\F_q$ 中可逆，约化映射
$R\to\F_q$ 与三平均相容。故每一步的所有坐标仍等于 $r$，不可能达到零。
\end{proof}

\begin{definition}\label{def:admissible}
固定 $n$，令
\[
 \PP_n=\{q:q\mid n,\ q\text{ 为素数},\ q\ne3\},\qquad d_n=|\PP_n|.
\]
称整数零和配置 $E$ 为\emph{算术合法}，如果对每个 $q\in\PP_n$，
其坐标模 $q$ 不全同余。
当 $\PP_n=\varnothing$ 时，该条件为空条件。
\end{definition}

算术合法性在下面的构造中直接施加于当前整数值，不在每一步重新取本原代表。
这使整个路径留在同一个离散格中。
对非零本原状态，它由引理 \ref{lem:G-divides-n} 与题目的 $G$ 条件完全等价。

\section{危险支持集与安全操作}\label{sec:safety}

本节除另有说明外均设 $n\ge11$。
以多重集记号 $a^r$ 表示 $r$ 个值等于 $a$ 的位置；它不表示数值乘方。
记 $f_E(a)$ 为实际值 $a$ 的重数，定义
\[
 \HH(E)=\{a:f_E(a)\ge3\}.
\]
我们称 $\HH(E)$ 的元素为\emph{三重值}。
一次整数三平均若输出仍算术合法，就称为\emph{安全操作}。
仅平均三个已经相等的值称为恒等操作。

\begin{lemma}[危险集]\label{lem:danger}
设 $E$ 算术合法，$q\in\PP_n$。若某个模 $q$ 余类含至少 $n-3$ 个位置，
则该余类唯一；记其补集为 $F_q$，有 $2\le|F_q|\le3$。
若不存在这样的余类，则删除任意三个位置后，余下位置模 $q$ 仍不全同余。
若存在，则任何不包含整个 $F_q$ 的三元组也具有这一性质。
\end{lemma}
\begin{proof}
由 $n-3>n/2$ 得唯一性。补集非空，否则违反合法性。
若补集只有一个位置，记主类余数为 $c$，例外余数为 $c+t$。
总和模 $q$ 等于 $nc+t=t$，故 $t=0$，矛盾。
其余结论来自补集的定义；删除三个位置后主类仍至少有 $n-6>0$ 个位置。
\end{proof}

\begin{lemma}[同时保护池]\label{lem:pool}
设一个位置池 $S$ 的所有值同余模 $3$，其中每个实际值最多出现两次。
则有以下结论。
\begin{enumerate}
\item 若 $|S|\ge d_n+3$，可在 $S$ 中选取安全、非恒等的整数三平均。
\item 若 $|S|\ge d_n+4$，对预先指定的数值 $a$，还可要求新均值不等于 $a$。
\end{enumerate}
\end{lemma}
\begin{proof}
逐个处理存在 $F_q$ 的素数 $q$。若当前可用池仍包含整个 $F_q$，留出其中一个位置；
否则无需动作。至多留出 $d_n$ 个位置，余下任意三元组都不包含任何完整 $F_q$，
因而由引理 \ref{lem:danger} 安全。其平均为整数；重数上界保证它非恒等。

第二种情形在保护后任取四个位置。若其每个三元组的均值均为 $a$，
比较只替换一个位置的两个三元组，便得四个值全相等，与重数至多二矛盾。
\end{proof}

\begin{lemma}[统一容量]\label{lem:capacity}
对所有 $n\ge11$，有 $n\ge d_n+10$。
\end{lemma}
\begin{proof}
$d_n=0,1$ 时显然。若 $d_n=2$，则 $n\ge14$：$11,12,13$ 都只有至多一个非三素因子。
若 $d_n\ge3$，除可能出现的 $2$ 外，其余非三素因子至少为 $5$，故
$n\ge2\cdot5^{d_n-1}\ge d_n+10$；最后一个不等式从 $d_n=3$ 起直接归纳。
\end{proof}

\begin{lemma}[重复值与仿射见证]\label{lem:repeated-safety}
以下操作只要平均为整数，均安全。
\begin{enumerate}
\item $(a,a,x)$，其中 $f_E(a)\ge3$。
\item $(a,b,x)$，其中 $a\ne b$ 且 $f_E(a),f_E(b)\ge2$。
\item 未选位置保留 $a,b$，而被选的至少两个值都是 $a,b$ 的整数仿射组合。
\end{enumerate}
第三项中的整数仿射组合是指 $ra+sb$，其中 $r,s\in\Z$、$r+s=1$。
\end{lemma}
\begin{proof}
固定 $q\in\PP_n$，反设输出全同余于 $c$。
第一项留下 $a$，故 $a\equiv c$，再由新均值同余于 $c$ 得 $x\equiv c$。
第二项留下 $a,b$，故二者都同余于 $c$，新均值再迫使 $x\equiv c$。
第三项中保留的 $a,b$ 同余于 $c$，因此至少两个被选值也同余于 $c$；
剩下一个由三元组之和被迫同余于 $c$。
每种情形都推出原配置全同余，矛盾。
\end{proof}

\begin{remark}[模 $2$ 的边界]\label{rem:two}
不能无条件将第二类操作改成 $(a,x,x)$。
例如 $(3^3,0^2,1^6,-15)$ 是合法的十二元零和状态，
但平均 $(3,0,0)$ 后全部变为奇数。
下文只在 $x=3b-2a$ 且操作后仍保留 $a,b$ 时使用 $(a,x,x)$，
其安全性来自引理 \ref{lem:repeated-safety} 的第三项，不需要在模 $2$ 下除以 $2$。
\end{remark}

\begin{lemma}[两值排除]\label{lem:two-values}
算术合法的整数零和配置不可能只有两个实际值 $a,b$，且 $a\not\equiv b\pmod3$。
\end{lemma}
\begin{proof}
设 $a$ 出现 $r$ 次，$0<r<n$。零和给出 $r(a-b)=-nb$。
对每个 $q\in\PP_n$，合法性保证 $q\nmid a-b$；对 $q=3$ 也由假设成立。
因此 $\gcd(a-b,n)=1$，进而 $n\mid r$，矛盾。
\end{proof}

\section{双三重值不变量的初始化}\label{sec:initialization}

定义增强不变量
\begin{equation}\label{eq:invariant}
 \II_n(E):\quad E\in\Z^n,\quad\sum E=0,\quad
 E\text{ 算术合法},\quad |\HH(E)|\ge2.
\end{equation}

\begin{proposition}[初始化]\label{prop:initialization}
设 $X\ne0$ 满足定理 \ref{thm:uniform} 的假设。
从 $3X$ 出发，至多两次安全、非恒等的整数三平均即可建立 $\II_n$。
\end{proposition}
\begin{proof}
初始所有值均被 $3$ 整除。乘以 $3$ 不改变有关素数下的不全同余性。
若已有两个三重值，则不需操作。

若恰有一个三重值 $a$，分两种情形。
当 $f_E(a)\ge5$ 时，取任意不同值 $b$ 并平均 $(a,a,b)$。
操作安全，留下至少三份 $a$，并产生三个不同于 $a$ 的新均值。
当 $f_E(a)=3$ 或 $4$ 时，$a$ 之外至少有 $n-4\ge d_n+6$ 个位置，
每个实际值至多出现两次。由引理 \ref{lem:pool}，可在这个池中安全地产生不同于 $a$ 的均值。

若没有三重值，先用引理 \ref{lem:pool} 选一个安全三元组。
新均值 $a$ 是唯一可能的三重值，其重数在 $3$ 与 $5$ 之间。
$a$ 之外至少有 $n-5\ge d_n+5$ 个位置，均为尚未改变的初值，仍被 $3$ 整除，
且每个实际值最多出现两次。再次应用引理 \ref{lem:pool}，
在这些位置中产生不同于 $a$ 的均值，即得两个三重值。
第二步不使用新生的 $a$，因此无需增加整数格的精度。
\end{proof}

\section{不变量的保持}\label{sec:closure}

记显式终端族为
\begin{equation}\label{eq:terminal}
 T_n(a)=(a^3,(-a)^3,0^{n-6}),\qquad a\ne0.
\end{equation}
它可连续三次平均 $(a,-a,0)$ 归零。

\begin{proposition}[保持或终止]\label{prop:closure}
设 $n\ge11$，$E$ 满足 $\II_n$。
则 $E$ 要么已经属于 \eqref{eq:terminal}，要么存在一次安全、非恒等的整数三平均，
其输出仍满足 $\II_n$。
\end{proposition}

下面按三重值的数目及其模 $3$ 余类分布，证明命题 \ref{prop:closure}。

\subsection{至少三个三重值}

若 $|\HH(E)|\ge4$，模 $3$ 的抽屉原理给出两个不同的同余三重值 $a,b$。
平均 $(a,a,b)$ 为安全整数操作；至少两个未参与的三重值完整保留。

若 $|\HH(E)|=3$，且有两个三重值 $a,b$ 同余模 $3$，记第三个为 $c$。
两种安全操作 $(a,a,b)$、$(b,b,a)$ 的均值分别为
\[
 d_1=(2a+b)/3,\qquad d_2=(a+2b)/3.
\]
由于 $d_1\ne d_2$，至少一个均值不同于未参与的 $c$；该新值与 $c$ 保持不变量。

剩下三个三重值分占三个模 $3$ 余类。
若另有非三重值 $x$，取与它同余模 $3$ 的三重值 $a$，平均 $(a,a,x)$；
其他两个三重值不动。
若没有其他实际值，记三个值为 $a,b,c$，平均 $(a,b,c)$。
若至少两个重数不小于 $4$，它们各用去一份后仍是三重值。
否则重数只能是 $3,3,n-6$。设 $c$ 的重数为 $n-6$，新均值为 $d$。
若 $d\ne c$，操作后 $c$ 至少还有四份，$d$ 有三份。
若 $d=c$，则 $a+b=2c$；结合零和得 $nc=0$，故 $c=0,b=-a$，
原状态即为 $T_n(a)$。

\subsection{恰有两个同余模 \texorpdfstring{$3$}{3} 的三重值}

记两个三重值为 $a,b$。
若 $f_E(a)\ge4$，平均 $(b,b,a)$，留下至少三份 $a$，并产生不同于 $a$ 的三重值；
$f_E(b)\ge4$ 时对称处理。
故只需考虑 $f_E(a)=f_E(b)=3$。
其余池 $S$ 有 $n-6\ge d_n+4$ 个位置，每个实际值最多出现两次。

若 $S$ 中存在与 $a,b$ 同余模 $3$ 的 $x$，考虑 $(a,a,x)$ 与 $(b,b,x)$。
二者都安全。如果它们都不能保持双三重值，必须分别有
\[
 2a+x=3b,\qquad 2b+x=3a.
\]
相减得 $5(a-b)=0$，矛盾。

否则，$S$ 只使用另外两个模 $3$ 余类。
若只使用其中一类，直接在 $S$ 内应用引理 \ref{lem:pool}。
若两类 $S_1,S_2$ 都非空，对跨类位置对 $(c,e)$ 考虑
$(a,c,e)$ 与 $(b,c,e)$。
这两种操作若被引理 \ref{lem:danger} 阻塞，只可能某个 $F_q$ 恰为所选的两个位置
$\{c,e\}$：若 $a$ 或 $b$ 属于 $F_q$，它的全部三份副本都属于 $F_q$，
而三元组只选了一份，故不可能包含整个 $F_q$。
因此至多 $d_n$ 个跨类位置对被阻塞，而
\[
 |S_1||S_2|\ge |S|-1=n-7>d_n.
\]
选择一个未被阻塞的对，两种操作均安全。
它们若都不能保持双三重值，就有
$a+c+e=3b$ 与 $b+c+e=3a$，推出 $4(a-b)=0$，仍为矛盾。

\subsection{恰有两个不同余模 \texorpdfstring{$3$}{3} 的三重值}

仍记三重值为 $a,b$。
若存在不同于 $a$、但与 $a$ 同余模 $3$ 的值 $x$，先尝试 $(a,a,x)$。
该操作安全；它不能保持双三重值，仅可能因为
\begin{equation}\label{eq:affine-neighbor}
 f_E(a)\le4,\qquad x=3b-2a.
\end{equation}
若这样的 $x$ 至少有两份，则改用 $(a,x,x)$，新均值为 $2b-a$，与 $a,b$ 都不同。
操作后仍有 $a,b$ 的副本，故由仿射见证引理，即使在模 $2$ 下也安全。
因此，在尚未找到保持操作的前提下，$a$ 的余类中至多有一个额外单点
$x=3b-2a$；对称地，$b$ 的余类中至多有一个额外单点 $y=3a-2b$。
若某个单点存在，相应三重值的重数不超过 $4$。
其余位置都在第三个模 $3$ 余类，记为池 $Z$。

\paragraph{两个额外单点均存在。}
由 $f_E(a)+f_E(b)+2\le10<n$，池 $Z$ 非空。
任取 $z\in Z$，平均 $(x,y,z)$。三重值 $a,b$ 不动；仿射见证保证安全。
又 $x-y=5(b-a)\ne0$，故操作非恒等。

\paragraph{恰有一个额外单点。}
不妨只有 $x$ 存在。若 $Z$ 为空，安全操作 $(a,a,x)$ 将输出只有 $a,b$ 的两值状态，
且 $a\not\equiv b\pmod3$，与引理 \ref{lem:two-values} 矛盾。
因此 $Z$ 非空。对 $z\in Z$，操作 $(x,b,z)$ 为安全非恒等整数操作。
若 $f_E(b)\ge4$，它保留 $a,b$ 两个三重值。
若 $f_E(b)=3$，则 $|Z|\ge n-8\ge3$。
每个实际值在 $Z$ 中最多两份，故至少有两个不同值；可选
$z\ne5a-4b$，使 $(x+b+z)/3\ne a$，从而保留 $a$ 和新三重值。

\paragraph{没有额外单点。}
若 $Z$ 为空，引理 \ref{lem:two-values} 已给出矛盾。
若 $f_E(a)=f_E(b)=3$，则 $|Z|=n-6\ge d_n+4$，在 $Z$ 内使用保护池引理。
若两个重数都不小于 $4$，任意 $(a,b,z)$ 均安全，并保留二者的三重性。

最后，不妨 $f_E(a)=3$、$f_E(b)\ge4$。
只要有 $z\ne2b-a$，平均 $(a,b,z)$ 即保持 $b$ 和不同于它的新三重值。
若不存在这样的 $z$，则池 $Z$ 全部等于 $2b-a$，其大小 $r$ 只能为 $1$ 或 $2$。
零和关系为
\begin{equation}\label{eq:last-collision}
 (3-r)(a-b)+nb=0.
\end{equation}
所有实际值都是 $a,b,2b-a$。
合法性说明每个 $q\in\PP_n$ 都不整除 $a-b$；又 $3\nmid a-b$，
故 $\gcd(a-b,n)=1$。
式 \eqref{eq:last-collision} 因而推出 $n\mid3-r$，与 $3-r\in\{1,2\}$ 矛盾。
命题 \ref{prop:closure} 的所有情况至此完成。\hfill$\square$

\section{整数能量、终止性与一般维数定理}\label{sec:termination}

\begin{lemma}[能量下降]\label{lem:energy}
设整数三元组 $(a,b,c)$ 的均值 $m$ 也是整数。
若三者不全相等，则平均使平方能量严格下降，且下降量至少为 $2$：
\begin{equation}\label{eq:energy}
 a^2+b^2+c^2-3m^2
 =\frac{(a-b)^2+(b-c)^2+(c-a)^2}{3}
 =(a-m)^2+(b-m)^2+(c-m)^2\ge2.
\end{equation}
\end{lemma}
\begin{proof}
展开得恒等式。三个整数偏差之和为零且不全为零，所以至少有两个非零偏差。
\end{proof}

\begin{proof}[定理 \ref{thm:uniform} 的证明]
零输入显然可达，故设 $X\ne0$。
从 $3X$ 出发，用命题 \ref{prop:initialization} 建立 $\II_n$。
只要不在终端族 $T_n$ 中，就用命题 \ref{prop:closure} 继续一次安全整数平均。
整个过程的
$\Phi(E)=\sum_i e_i^2$ 是非负整数，并由引理 \ref{lem:energy} 严格下降。
因此不可能无限继续，必在有限步到达 $T_n$；随后三次整数平均归零。

该序列全部在原来的 $n$ 个位置上执行。由仿射相容性，在 $X$ 上执行同样的位置序列，
每个中间值恰为上述整数值的三分之一，故均在 $\frac13\Z$ 中。
所有操作都是非恒等整数平均，故从 $3X$ 出发的总操作数满足粗界
\begin{equation}\label{eq:length}
 L\le\frac92\sum_i x_i^2.
\end{equation}
最后，若原整数输入有公因子，先在证明中取本原代表，应用上述构造，
再把同一操作序列整体乘回该整数因子。此举不增加分母。
\end{proof}

\begin{remark}
式 \eqref{eq:length} 是关于数值高度的界。
若 $H=\max_i|x_i|$，它给出 $L=O(nH^2)$，而不是关于 $\log H$ 的多项式界。
安全性、不变量保持和终止性分别由前面的全称引理承担；有限数值测试不用于替代它们。
\end{remark}

\section{七至十元的独立证明}\label{sec:small}

本节只用于定理 \ref{thm:main} 中尚未被定理 \ref{thm:uniform} 覆盖的四个维数。
首先给出一个适用于小维数入口的高重数引理。

\subsection{安全下降至高重数核心}

本小节允许 $n\ge7$，并取非零本原零和状态 $X$，满足 $G(X)=3^a$。
对每个存在模 $q$ 主类的非三素因子，仍用 $F_q$ 表示其例外集；
记这样的危险素数数量为 $D$。
称三元组为\emph{补集安全}，若未选位置对每个非三素因子仍不全同余。

\begin{lemma}[高重数核心]\label{lem:core}
若没有补集安全的非恒等整数三平均，则某个实际值至少出现 $n-D-2$ 次。
\end{lemma}
\begin{proof}
危险集仍唯一且大小为 $2$ 或 $3$；此时 $n-3>n/2$ 对 $n\ge7$ 仍成立。
令 $U=\bigcup_qF_q$、$W=\{1,\ldots,n\}\setminus U$。
有 $|W|\ge n-3D\ge4$：$D=0,1$ 时由 $n\ge7$ 得到，
$D\ge2$ 时使用 $n\ge2\cdot5^{D-1}\ge3D+4$。
任意包含两个 $W$ 位置的三元组都不能包含完整危险集，因而只要整平均非恒等，就补集安全。

$W$ 不能同时出现三个模 $3$ 余类。若其中两个余类各有至少两个位置，
则全局第三余类必须为空；否则从这两个 $W$ 余类各取一个位置，再加第三余类的位置，便得安全操作。
若两个全局余类也各至少有三个位置，用其中两个 $W$ 位置与本类任意第三位置组队，
可知每个全局余类均为实际常值。这产生两个不同余模 $3$ 的实际值，
由引理 \ref{lem:two-values} 的同一算术证明不可能。
因此其中一个全局余类至多两个位置，另一类为至少 $n-2$ 重的常值，结论成立。

剩下 $W$ 的一个模 $3$ 余类至少占 $|W|-1\ge3$ 个位置。
它们必须等于同一实际值 $u$；任意与 $u$ 同余模 $3$ 的全局位置配上这两个 $W$ 位置，
同样必须等于 $u$。所有非 $u$ 位置因此只在另外两个余类中，且所有 $F_q$ 都在这些例外位置内。

若两个例外余类的大小分别为 $r_1,r_2>0$，任取一对跨类例外，配一个 $W$ 中的 $u$，
即为非恒等整数三元组。它被阻塞只能因为该例外对恰为某个 $F_q$。
故 $r_1r_2\le D$，从而 $r_1+r_2\le D+1$。

若例外只有一个余类，记其大小为 $r$。
对每个 $F_q$ 任取两个位置连一条边，得到 $r$ 个顶点、至多 $D$ 条边的图。
若 $r\ge D+3$，该图至少有三个连通分量。
从三个不同分量各取一个顶点，所得三元组不包含任何完整 $F_q$，且平均为整数，
所以其三个实际值必须相等。
固定不同分量中的代表，再逐个替换同一分量的代表，可推出全部例外实际值相等。
于是又得到两个不同余模 $3$ 的实际值，矛盾。
故 $r\le D+2$，即得所求重数界。
\end{proof}

不断执行补集安全的非恒等整数平均，并在每一步后取本原整数代表。
补集保留的余数差保证约去的公因子不含有关素数，故 $G$ 条件保留；
平方能量在平均及随后约分时严格下降。因此有限步后到达零或引理 \ref{lem:core} 的高重数核心。

定义标准族
\begin{equation}\label{eq:B}
 \B_n(u,v)=\bigl(u^{n-4},v^3,-(n-4)u-3v\bigr).
\end{equation}
对 $\gcd(u,v)=1$ 的整数参数，直接计算得
\begin{equation}\label{eq:B-G}
 G(\B_n(u,v))=\gcd(u-v,n).
\end{equation}
在 $n=7,8$ 时，非三素因子只有一个，故核心至少有 $n-3$ 个相等值 $u$。
从三个例外中保留一个模该素数不同于 $u$ 的值 $t$，平均另外两个例外和一个 $u$，
就得到合法的 $\B_n$。合法性由
$u-v\equiv(t-u)/3\pmod q$ 得到。
十元核心至少有六重值，其入口在第 \ref{subsec:ten-bridge} 小节处理。

\subsection{七元：精确有限下降证书}\label{subsec:seven}

在 $\B_7(u,v)$ 中令 $x=u,y=u+v$。
对本原整数参数有
\[
 Q_7(x,y)=x^2-xy+2y^2,\qquad
 \Phi(\B_7)=6Q_7,\qquad
 Q_7\equiv2(2x-y)^2\pmod7.
\]
因此合法性恰为 $7\nmid Q_7$。
平均单点 $-3y$ 与两个 $v$，或与两个 $u$，分别给出射影矩阵
\begin{equation}\label{eq:seven-generators}
 A=\mat{3}{0}{1}{-1},\qquad B=\mat{-3}{3}{-1}{0}.
\end{equation}
其真实作用均为所列矩阵的 $1/3$ 倍，返回时按新的三重值重新记 $u,v$。
两者行列式只含素因子 $3$，且模 $7$ 分别将 $2x-y$ 乘以 $6$ 和 $1$，故保持合法性。

\begin{lemma}[七元有限证书]\label{lem:seven-certificate}
设 $(x,y)$ 为本原整数对，$7\nmid Q_7(x,y)$ 且 $y\ne0$。
附录 \ref{app:seven} 所列十八个 $A,B$ 词中，至少一个严格降低输出本原代表的 $Q_7$。
\end{lemma}

引理的证明是附录所给的完整有限代数验证：模 $3^7$ 的 $2916$ 个本原射影类
压缩为五种赋值签名，每种签名对应一组显式整系数二次多项式。
精确实根隔离证明其正区间覆盖所有合法非终端有理方向。
这里检验的是有限代数证书，不是对整数高度的截断搜索。

由该引理，正整数 $Q_7$ 有限次下降后到达 $y=0$。
此时状态为 $(u^3,(-u)^3,0)$，三次 $(u,-u,0)$ 即归零。
结合高重数入口，七元充分性成立。

\subsection{八元：三进欧几里得下降}\label{subsec:eight}

仍令 $x=u,y=u+v$。有
\[
 Q_8(x,y)=2x^2+3y^2,\qquad \Phi(\B_8)=4Q_8,
\]
而合法本原参数恰为 $y$ 是奇数。
附录 \ref{app:eight} 列出三个真实平均词，其射影作用为
\begin{equation}\label{eq:eight-macros}
 R:(x,y)\mapsto(x+y,-y),\quad
 S:(x,y)\mapsto(y-x,y),\quad
 F:(x,y)\mapsto(3x+3y,2x-4y).
\end{equation}
前两者的合成提供 $x\mapsto x\pm2y$，所以任意整数倍的这种平移均可执行。

若 $|y|>1$，先在必要时使用 $S$，使 $x$ 也为奇数。
利用平移选择 $x_0\equiv x\pmod{2y}$ 满足
\[
 |x_0-2y|<|y|.
\]
最近剩余总能达到不严格的不等式；若等号成立，则 $y\mid x$，
与本原性和 $|y|>1$ 矛盾。
施行 $F$ 得到 $(3(x_0+y),2(x_0-2y))$。
两坐标均为偶数，而除以 $2$ 后第二坐标为奇数。
取本原代表后仍有奇数 $y'$，并且 $|y'|\le|x_0-2y|<|y|$。
因此正整数 $|y|$ 严格下降。

到 $|y|=1$ 时，整体改符号使 $y=1$，平移将 $x$ 化为 $0$ 或 $1$；
后一种再用 $S$ 即化为 $0$。
代表 $\B_8(0,1)=(0^4,1^3,-3)$ 的四步收尾为：
先平均 $(-3,0,0)$，再三次平均 $(-1,0,1)$。
八元充分性得证。

\subsection{九元：固定三进网络}\label{subsec:nine}

把九个位置排成 $3\times3$ 方阵，先平均每一行，再平均每一列。
六次操作后全部等于总体均值。
更一般地，将 $3^k$ 个位置标为长度为 $k$ 的三进制串，逐坐标平均其余坐标固定的三个位置，
即可在 $k3^{k-1}$ 步后平均整个块。

\subsection{十元：六重核心的入口}\label{subsec:ten-bridge}

由引理 \ref{lem:core}，只需处理
$E=(u^6,a,b,c,d)$。令四个例外差为 $e_t=t-u$，则
\[
 e_a+e_b+e_c+e_d=-10u.
\]
合法性保证至少一个 $e_t$ 为奇数，且至少一个不被 $5$ 整除。
若某个差与 $10$ 互素，保留对应例外 $t$，平均其他三个例外，即一步进入合法 $\B_{10}$。

否则，每个奇差都被 $5$ 整除，每个非五倍数差都为偶数。
零和模 $2,5$ 迫使四差模 $10$ 为两个 $5$ 与两个互为相反数的非零偶数。
因此可重新编号，使 $e_a+e_b$ 与 $10$ 互素。
置
\[
 q=\frac{-5u-a-b}{3},\qquad
 r=\frac{-5u+2a+2b}{9},\qquad
 s=\frac{-11u-a-b}{27}.
\]
依次平均 $(u,c,d)$、$(q,a,b)$、$(u,q,r)$、$(u,q,r)$。
完整重数账本为
\[
 u^6abcd\longrightarrow u^5abq^3\longrightarrow u^5q^2r^3
 \longrightarrow u^4qr^2s^3\longrightarrow u^3rs^6.
\]
输出是 $\B_{10}(s,u)$，且
\[
 r=-6s-3u,\qquad s-u=-\frac{40u+e_a+e_b}{27}.
\]
故参数差在模 $2,5$ 下均非零，所得标准核合法。

\subsection{十元：局部与实数联合下降}\label{subsec:ten-descent}

在 $\B_{10}(u,v)$ 中，将 $(u,3u+2v)$ 取本原整数代表，记为 $(X,Y)$。
合法参数分成两种类型：
\begin{equation}\label{eq:ten-type}
 h(X,Y)=
 \begin{cases}
 1,&X,Y\text{ 都为奇数且 }X+Y\equiv0\pmod4,\\
 2,&X,Y\text{ 奇偶相反},
 \end{cases}
 \qquad 5\nmid Y.
\end{equation}
相应的本原 $(u,v)$ 分别是
$(X,(Y-3X)/2)$ 与 $(2X,Y-3X)$。
取正整数高度
\begin{equation}\label{eq:ten-height}
 W(X,Y)=h(X,Y)^2(5X^2+Y^2)=\frac13\Phi(\B_{10}(u,v)).
\end{equation}
附录 \ref{app:ten} 列出七个真实返回词，其射影矩阵依次为
\begin{equation}\label{eq:ten-matrices}
\begin{gathered}
 M_1=\mat{1}{1}{-5}{-1},\quad M_2=\mat{1}{0}{-5}{2},\quad
 M_3=\mat{1}{-2}{0}{3},\quad M_4=\mat{1}{1}{5}{-3},\\
 M_5=\mat{4}{1}{0}{-3},\quad M_6=\mat{1}{-2}{5}{2},\quad
 M_7=\mat{1}{0}{-10}{-1}.
\end{gathered}
\end{equation}

\begin{lemma}[十元有限证书]\label{lem:ten-certificate}
每个满足 \eqref{eq:ten-type} 的本原整数方向，若 $X\ne0$，则有一个上述真实返回词，
其输出仍合法并严格降低本原高度 $W$。
\end{lemma}

附录给出该引理所需的全部矩阵、位置词、模 $96$ 签名及精确实数判定过程。
行列式只含素因子 $2,3$，因此共同约分因子及输出模 $4$ 类型由有限模数决定。
$5120$ 个本原局部类仅产生十二种签名。
对应二次不等式的精确验证说明：唯一相容的未覆盖射影方向是 $X=0$。

因 $W$ 为正整数，反复应用引理必有限步到达 $X=0$，即 $u=0$。
此时保留一个零位置，对其余九个零和位置使用第 \ref{subsec:nine} 小节的网络即可。
十元充分性得证。

\begin{proof}[定理 \ref{thm:main} 的证明]
必要性是命题 \ref{prop:necessary}。
对 $n\ge11$，引理 \ref{lem:G-divides-n} 将条件转为算术合法性，
定理 \ref{thm:uniform} 给出充分性。
$n=7,8,10$ 分别由本节的独立构造证明，$n=9$ 由固定网络证明。
最后用引理 \ref{lem:affine-equivariance} 恢复原有理输入和原均值。
\end{proof}

\section{四、五、六元的完整判定}\label{sec:boundary}

本节始终将输入中心化并清除分母。
四元可达性有直接的充要条件；五元有确定性递归和位长多项式复杂度；
六元的可判定性也已解决，但下面仅给出伪多项式算法。
记 $H$ 为所使用的整数零和输入的最大坐标绝对值。

\subsection{成功路径不产生新分母}

\begin{proposition}[小维数不产生新分母]\label{prop:small-integral}
设 $3\le n\le6$。从整数零和状态出发，任何最终归零的三平均路径都始终为整数。
\end{proposition}
\begin{proof}
假设第一次非整数操作产生三个相等值 $c=t/3$，其中 $3\nmid t$；其他位置仍为整数。
此后追踪如下性质：最低三进赋值为 $-r<0$，恰有三个位置 $T$ 达到该赋值，且它们实际相等；
其他位置的赋值至少为 $-r+1$。
在每个这样的状态中，重新记 $t=3^r c$，其中 $c$ 是这三个位置的共同值，故 $t$ 为三进单位。
若下次选取的三元组与 $T$ 相交于一个或两个位置，乘以 $3^r$ 后其和模 $3$ 为
$t$ 或 $2t$，非零。平均后最低赋值下降一，并恰在新的三个相等位置取得，性质继续保持。
若恰选 $T$，操作恒等。
若与 $T$ 不交，只有 $n=6$ 才可能发生；由零和，输出为
$(c^3,(-c)^3)$。整体除以 $c$ 后是 $(1^3,(-1)^3)$，其所有值模 $2$ 为 $1$，
故由同余障碍不能归零。
因此第一次产生非整数后不可能成功，矛盾。
\end{proof}

\begin{example}[判据在五、六元失效]
五元状态 $(-3,0,1,1,1)$ 满足 $G=1$，但没有非恒等整数三平均，因而由命题
\ref{prop:small-integral} 不可达。
六元状态 $(-3,0,0,1,1,1)$ 也满足 $G=1$；其唯一非恒等整数操作产生
$(-1^3,1^3)$ 的模 $2$ 障碍，因此不可达。
这说明定理 \ref{thm:main} 的统一维数阈值 $7$ 不能降低。
\end{example}

\subsection{四元的显式充要条件}

\begin{proposition}[四元判据]\label{prop:four}
四个有理数可经有限次三平均变为常值，当且仅当至少一个初始数已等于总体平均值。
\end{proposition}
\begin{proof}
中心化后，只需证明可达性等价于初始已有零坐标。
若有零坐标，保留它并平均其余三个即可。
反之，一次操作若留下数值 $a$，所选三个的均值为 $-a/3$。
因此全非零的零和状态操作后仍全非零，永远不能达到零。
该论证直接适用于有理数，无需整数化。
\end{proof}

三个位置则总可一步平均；这与四元判据一起处理了 $n=3,4$。

\subsection{五元的闭合二维族与陷阱}

第一次三平均之后，五元零和状态必可写成
\begin{equation}\label{eq:five-family}
 E=(a,b,c,c,c),\qquad a+b+3c=0.
\end{equation}
令 $u=a-c,v=b-c$。非零状态的有理对 $(u,v)$ 可以整体缩放为本原整数对。
由于 $5c+u+v=0$，这两个差值在整体尺度之外完全确定状态。
重新以新生三重值为基准，三种操作给出表 \ref{tab:five-moves} 的射影变换。
两个剩余单点的顺序可以交换。

\begin{table}[htbp]
\centering
\caption{五元闭族的全部操作类型}\label{tab:five-moves}
\begin{tabular}{cc}
\toprule 选取类型 & 新差值对的射影代表\\\midrule
$(a,c,c)$ & $(-u+3v,-u)$\\
$(b,c,c)$ & $(3u-v,-v)$\\
$(a,b,c)$ & $(-u-v,-u-v)$\\
$(c,c,c)$ & $(u,v)$\\
\bottomrule
\end{tabular}
\end{table}

考虑本原整数对组成的集合
\begin{equation}\label{eq:five-trap}
 \mathcal T=\{(u,v):u\equiv v\not\equiv0\pmod3\}.
\end{equation}
它是所有操作下的闭集。事实上，前三种原始输出分别同余于
$(-u,-u)$、$(-v,-v)$、$(-u-v,-u-v)$，两坐标均为相同的非零余数；
共同约分因子不含 $3$，约分后性质仍成立。
而成功终端 $u+v=0$ 不在该集合中。因此 $\mathcal T$ 中的状态均不可达。

\subsection{五元递归的充要性与终止性}

定义布尔判据 $\mathscr R(u,v)$。零对直接接受；非零有理对先取本原整数代表，
再按以下规则计算：
\begin{equation}\label{eq:five-recursion}
\mathscr R(u,v)=
\begin{cases}
 \mathrm{true},&u+v=0,\\
 \mathscr R(v-u/3,-u/3),&u\ne0,\ 3\mid u,\ 3\nmid v,\\
 \mathscr R(u-v/3,-v/3),&v\ne0,\ 3\mid v,\ 3\nmid u,\\
 \mathrm{false},&\text{其他情形}.
\end{cases}
\end{equation}
两个递归分支不能同时适用，因为输入对本原。

\begin{theorem}[五元递归判据]\label{thm:five}
闭族 \eqref{eq:five-family} 可达零，当且仅当 $\mathscr R(a-c,b-c)$ 为真。
对任意五元输入，枚举第一次选择的十个三元组，再对所得闭族应用该判据，即给出完整判定。
接受时可恢复一条真实原位置操作序列。
\end{theorem}
\begin{proof}
若 $u+v=0$，则 $5c+u+v=0$ 给出 $c=0$；平均 $(a,b,c)$ 即全部归零。
若 $u+v\ne0$，选 $(a,b,c)$ 会产生两个相等的非零差值，
取本原代表后落入 $\mathcal T$，所以它不可能是成功分支。

若 $u,v$ 都不是三倍数，则前两种操作也立即落入 $\mathcal T$。
若恰有一个坐标为零，则与它对应的 $(a,c,c)$ 或 $(b,c,c)$ 只是恒等操作；
其他非恒等操作仍落入陷阱，故应拒绝。
其余情形恰有一个非零坐标被 $3$ 整除。
若 $u=3q\ne0$，选 $(a,c,c)$ 后除去射影尺度 $3$，新对为 $(v-q,-q)$。
它仍本原，因为 $\gcd(v-q,q)=\gcd(v,q)=1$。
另一个非恒等选择 $(b,c,c)$ 则落入陷阱。
对 $v$ 是三倍数的情形同理，得到 \eqref{eq:five-recursion} 的第三行。
所以除恒等操作外，每个非终端状态至多有一个可能成功的分支。

每个递归步都严格降低整数 $|u|+|v|$。例如 $u=3q\ne0$ 时，
\begin{equation}\label{eq:five-l1}
 |v-q|+|q|\le |v|+2|q|<|v|+3|q|=|v|+|u|.
\end{equation}
递归因此必定终止，且根据前述分支排除，其真假与可达性等价。
一般五元输入若非零，任何成功路径都有第一步；枚举全部十个选择是完备的。
每次接受的递归步都指定表 \ref{tab:five-moves} 的一个实际三元组，
最后接终端操作，因此也给出操作证书。
\end{proof}

\subsection{五元的平方对数步数界}

\begin{theorem}[五元复杂度]\label{thm:five-complexity}
设五元整数零和输入的高度为 $H$。
上述判定在 $O((\log(2+H))^2)$ 次大整数递归步骤内完成，
所处理整数均只有 $O(\log(2+H))$ 位。
因此五元判定及成功路径构造均具有输入位长多项式复杂度。
\end{theorem}
\begin{proof}
先考虑一个本原非终端对。将其中非零的三倍数称为活动坐标，
必要时交换两个坐标，使其在第一位。
写活动坐标为 $a=3^kt$，其中 $k\ge1$、$3\nmid t$，另一个坐标为 $b$，故 $3\nmid b$。
从这一步起连续消耗活动坐标的三进因子，称为一段\emph{级联}。
第一步后活动坐标位于第二位，此后继续使用 \eqref{eq:five-recursion} 的第三行，
直到该坐标成为三进单位。
在前 $k-1$ 步，两坐标恰有一个为三倍数，因而不可能提前接受或拒绝。
直接归纳得到整段 $k$ 步的输出
\begin{equation}\label{eq:five-cascade}
 C_k(a,b)=\left(b-d_kt,(-1)^kt\right),\qquad
 d_k=\frac{3^k+(-1)^{k-1}}4\in\Z.
\end{equation}
段末第二坐标为三进单位；若递归继续，第一坐标就是下一段的活动坐标。

令
\[
 \alpha_k=\frac{d_k}{3^k}
 =\frac14+\frac{(-1)^{k-1}}{4\cdot3^k},\qquad
 \beta_k=\frac{(-1)^k}{3^k}.
\]
则 $C_k(a,b)=(b-\alpha_ka,\beta_ka)$，且
$0<\alpha_k\le1/3$、$|\beta_k|\le1/3$。
对于两段连续级联，记最终输出为 $(a_2,b_2)$，有
\begin{equation}\label{eq:five-contraction}
 |a_2|\le\frac49|a|+\frac13|b|,\qquad
 |b_2|\le\frac19|a|+\frac13|b|.
\end{equation}
因此以 $H_0=\max(|a|,|b|)$ 记初始差值高度，
每两段级联后高度至多为之前的 $7/9$；单段至多放大为 $4/3$ 倍。
非零整数对的高度至少为 $1$，故级联总数为 $O(\log(2+H_0))$，
每个级联起点高度均不超过 $4H_0/3$。
又 $3^k\le|a|$，所以每段的步数 $k$ 也为 $O(\log(2+H_0))$。
总递归步数遂为平方对数。
式 \eqref{eq:five-l1} 还保证所有中间坐标绝对值不超过初始 $|a|+|b|$，
因此位长保持对数级。

对一般五元输入，第一次三平均的均值记为 $c$，留下 $a,b$。
用 $(3a-3c,3b-3c)$ 清除差值分母；两坐标绝对值均不超过 $6H$。
取本原代表只会减小高度。首次选择只有十种，故上述复杂度界仍成立。
有理输入的中心化和清分母具有多项式位复杂度，结论随之成立。
\end{proof}

这里的平方对数是递归步骤数，不是对每一步大整数运算都按单位位成本计费。
将该界进一步改进为 $O(\log(2+H))$ 并未在本文中证明，也不影响五元完整可判定性。

\subsection{六元的完整有限图算法}

\begin{theorem}[六元判定]\label{thm:six}
对整数零和六元输入，存在至多 $O(H^3)$ 个状态的有限图，
使输入可达零当且仅当该图中存在通往零的路径。
判定成功时可恢复实际三平均序列，失败时可穷尽全部可能的整数后继状态。
\end{theorem}
\begin{proof}
由命题 \ref{prop:small-integral}，任何成功路径均不产生非整数，所以只需保留整数平均。
平均保持初值的凸包，故坐标始终属于 $[-H,H]\cap\Z$。
第一次操作后，每个状态都可表示为
\[
 (a,b,d,c,c,c),\qquad c=-\frac{a+b+d}{3}.
\]
把配置视为多重集，枚举 $a,b,d$ 的有序取值已覆盖全部这样的状态，
数量至多 $(2H+1)^3$。再加入原始状态。
每个状态只有 $\binom63=20$ 个位置三元组；保留平均为整数且非恒等的边。
能量严格下降，因此不存在有向环。
在这个有限图中作完整可达性判定即可；命题 \ref{prop:small-integral} 保证没有丢弃任何成功路径。
记录前驱位置或数值重数，便可重建真实操作。
\end{proof}

上述状态数界另乘每次整数运算的对数位成本，仍属于关于 $H$ 的伪多项式界。
本文没有给出六元的简洁静态算术充要条件，也没有证明它具有输入位长多项式算法。
这两个问题的保留不等于六元可判定性尚未解决。

\begin{table}[htbp]
\centering\small
\caption{各维数的完整判定及复杂度边界}\label{tab:dimension-summary}
\begin{tabular}{cp{0.43\textwidth}p{0.33\textwidth}}
\toprule 维数 & 判定依据 & 已证复杂度\\\midrule
$3$ & 一步平均全部位置 & 常数次操作\\
$4$ & 某个初值等于总体均值 & 位长多项式\\
$5$ & 十个首步选择与确定性递归 & $O((\log(2+H))^2)$ 个递归步；位长多项式\\
$6$ & 完整整数状态图 & $O(H^3)$ 个状态；伪多项式\\
$n\ge7$ & 本原中心化差分 gcd 为三幂 & 判定为位长多项式；统一构造长度的位长多项式界未证\\
\bottomrule
\end{tabular}
\end{table}

\section{算法含义与证据范围}\label{sec:algorithm}

对 $n\ge7$，判定只需中心化、清除分母、计算坐标 gcd 和差分 gcd，再不断除去因子 $3$。
有理数分母最小公倍数的位长不超过各输入分母位长之和；整数加减乘除及 gcd 都具有多项式位复杂度。
故定理 \ref{thm:main} 给出关于输入总位长的多项式时间判定算法。

$n\ge11$ 的构造每步只需检查实际值、重数和有限个余数。
保护池至多留出每个危险素数的一份见证，之后只查看三或四个位置。
主体证明既不寻找矩阵群生成元，也不要求解一个预先给定的运输矩阵分解。
允许依赖当前状态选择三元组，是构造的组成部分。

本文证据分为两类：第 \ref{sec:safety}--\ref{sec:termination} 节是对任意维数与任意整数大小的直接证明；
七、十元的有限覆盖则是计算机辅助的精确代数证书。
后者的有限模空间、矩阵、真实位置词及多项式定义全部列在附录，补充材料提供仅用整数和有理数的验证程序。
全文没有使用浮点抽样来证明无限区间覆盖。

全维数执行器另做有界状态、较大整数坐标和带标签完整路径重放。
这些重放用于防止账本和实现错误，并不承担全称量词。
其中模 $2$ 的反例与仿射修复单独检查，以避免把只适用于奇素数的除法规则用于偶维数。

当前尚未由本文证明的是：七、八、十元的一位额外三进精度界，以及所有维数统一的位长多项式构造长度。
此外，六元的位长多项式判定与五元递归的线性对数步数改进也不在本文结论之内。
这些问题与四、五元及 $n\ge7$ 已有的多项式时间判定应予区别。

\appendix
\section{七元证书的完整有限数据}\label{app:seven}

用字母 $0,1$ 分别表示 \eqref{eq:seven-generators} 的 $A,B$。
字串按时间从左向右执行，故 $01$ 的矩阵为 $BA$。
依次使用十八个词
\begin{equation}\label{eq:seven-words}
\begin{gathered}
0,\ 1,\ 01,\ 11,\ 001,\ 010,\ 101,\ 110,\\
0100,\ 0111,\ 1100,\ 1111,\ 00111,\ 01110,\ 10111,\ 11110,\\
0111101,\ 1111101.
\end{gathered}
\end{equation}
记词矩阵为 $M_w$，取
\[
 e_w(z)=\min\{v_3((M_wz)_1),v_3((M_wz)_2)\}.
\]
由于 $\det M_w=\pm3^{|w|}$ 且输入本原，共同因子只有 $3$；
实际本原输出恰为 $3^{-e_w}M_wz$，整体符号无关。

模 $3^7$ 的本原射影类有完整代表集
\[
 \{(1,t):0\le t<3^7\}\ \cup\
 \{(3t,1):0\le t<3^6\}.
\]
它们分成表 \ref{tab:seven-types} 的五类。各类的十八项 $e_w$ 签名为
\begin{align*}
\sigma_1={}&(0,0,1,1,1,1,1,1,1,2,1,2,2,3,2,3,3,3),\\
\sigma_2={}&(1,0,1,1,2,1,1,1,1,2,1,2,3,2,2,2,4,3),\\
\sigma_3={}&(0,0,1,1,1,2,1,2,3,2,3,2,2,2,2,2,3,3),\\
\sigma_4={}&(0,0,1,1,1,2,1,2,2,2,2,2,2,2,2,2,3,3),\\
\sigma_5={}&(0,1,1,1,1,1,2,1,1,2,1,2,2,2,3,2,3,4).
\end{align*}
所有赋值均小于 $7$，故模 $3^7$ 足以精确确定它们。

\begin{table}[htbp]
\centering\small
\caption{七元的五种局部签名与实射影未覆盖点}\label{tab:seven-types}
\begin{tabular}{clrrl}
\toprule
类 & 局部条件 & 类数 & 未覆盖有限斜率 & 相容未覆盖斜率\\
\midrule
1 & $3\nmid x,\ y/x\equiv0\pmod3$ &729& $0,2$ & $0$\\
2 & $3\nmid x,\ y/x\equiv1\pmod3$ &729& 无 & 无\\
3 & $3\nmid x,\ y/x\equiv2\pmod9$ &243& $0,2$ & $2$\\
4 & $3\nmid x,\ y/x\equiv5,8\pmod9$ &486& $0,2$ & 无\\
5 & $3\mid x,\ 3\nmid y$ &729& 无 & 无\\
\bottomrule
\end{tabular}
\end{table}

令 $K=\mat{2}{-1}{-1}{4}$，则 $z^TKz=2Q_7(z)$。
对每个签名和每个词，生成齐次二次型
\begin{equation}\label{eq:seven-D}
 D_w(z)=z^T\bigl(3^{2e_w}K-M_w^TKM_w\bigr)z.
\end{equation}
$D_w(z)>0$ 正是该词的严格下降条件。
按附录 \ref{app:exact-roots} 的有限算法检验 $D_w(1,t)$，得到表中的未覆盖点；
无穷方向在五类中都被覆盖。
斜率 $0$ 是终端，斜率 $2$ 满足 $Q_7(x,2x)=7x^2$，不合法。
这证明引理 \ref{lem:seven-certificate}。

\section{八元宏的真实位置账本}\label{app:eight}

令 $U=u,V=v,W=-4u-3v$。下表的每一项表示一次对三个当前实际值的平均；
出现相同数值两次时取两个不同位置。
\begin{table}[htbp]
\centering
\caption{实现 \eqref{eq:eight-macros} 的原子词}\label{tab:eight-words}
\begin{tabular}{c>{\raggedright\arraybackslash}p{0.79\textwidth}}
\toprule 宏 & 按时间排列的三元组\\\midrule
$R$ & $(V,V,U)$；$(V,U,U)$；
$\bigl(W,(U+2V)/3,(2U+V)/3\bigr)$；
$\bigl((U+2V)/3,(2U+V)/3,U\bigr)$。\\[0.6em]
$S$ & $(W,V,U)$；$(-U-2V/3,-U-2V/3,U)$；
$(-U-2V/3,U,U)$；
$(-U/3-4V/9,V,U/3-2V/9)$ 连做两次；
$(V/9,V/9,U/3-2V/9)$。\\[0.6em]
$F$ & $(V,U,U)$；$\bigl(W,(2U+V)/3,(2U+V)/3\bigr)$；$(V,U,U)$。\\
\bottomrule
\end{tabular}
\end{table}
逐次替换可得真实的 $(u,v)$ 输出矩阵分别为
\[
 R_{uv}=\mat{2/3}{1/3}{-1}{-2/3},\qquad
 S_{uv}=\mat{0}{1/9}{1/9}{0},\qquad
 F_{uv}=\mat{2/3}{1/3}{-8/9}{-7/9}.
\]
在 $(x,y)=(u,u+v)$ 下，分别为 \eqref{eq:eight-macros} 所列整数矩阵的
$1/3,1/9,1/9$ 倍。形式独立的 $u,v$ 下位置重数充足；
数值特化造成重合时，重数相加，不会破坏任一次选取的可执行性。

\section{十元证书的完整有限数据}\label{app:ten}

\subsection{七个真实词与实际尺度}

置
\begin{gather*}
 U=u,\qquad V=v,\qquad W=-6u-3v,\\
 q=\frac{-5u-2v}{3},\quad r=\frac{u+2v}{3},\quad
 s=\frac{2u+v}{3},\quad t=-2u-\frac v3.
\end{gather*}
七个矩阵 $M_i$ 由表 \ref{tab:ten-words} 的真实平均词实现。
表中的 $\lambda_i$ 满足：在未取本原代表的坐标 $(u,3u+2v)$ 中，真实输出为
$\lambda_iM_i(u,3u+2v)^T$。

\begin{table}[htbp]
\centering\small
\caption{十元矩阵的原子词、行列式与实际尺度}\label{tab:ten-words}
\begin{tabular}{crrp{0.60\textwidth}}
\toprule
$i$ & $\det M_i$ & $\lambda_i$ & 依次平均的当前值\\
\midrule
1&4&$1/6$&$(W,V,U),(V,U,U),(V,U,U)$\\
2&2&$-1/9$&$(W,V,U),(q,U,U),(V,V,U),(q,r,U),(q,r,U)$\\
3&3&$-1/9$&$(W,U,U),(V,U,U),(V,s,U),(V,s,U)$\\
4&$-8$&$1/6$&$(W,U,U),(V,U,U),(V,U,U)$\\
5&$-12$&$1/9$&$(W,V,U),(V,V,U),(r,U,U),(r,U,U)$\\
6&12&$-1/9$&$(V,U,U),(W,s,U),(V,s,U),(V,s,U)$\\
7&$-1$&$-1/9$&$(W,V,V),(V,U,U),(t,s,U),(t,s,U),(s,U,U)$\\
\bottomrule
\end{tabular}
\end{table}

\subsection{局部签名}

对本原方向 $z=(X,Y)$，定义
\[
 g_i(z)=\gcd\bigl((M_iz)_1,(M_iz)_2\bigr)>0.
\]
由伴随矩阵恒等式，$g_i\mid|\det M_i|$，故它只含素因子 $2,3$。
记约分后的输出类型为 $h_i'$；若不满足 \eqref{eq:ten-type} 的奇偶条件，则记为 ``--''。
模数 $96=2^5\cdot3$ 足以决定 $g_i$ 及输出模 $4$ 的类型：最大二赋值为 $3$，
除后还需模 $4$，故模 $2^5$ 足够；最大三赋值为 $1$。
约去 $3$ 因子时，它在模 $4$ 下可逆。

在所有模 $96$ 对中，只保留模 $2$、模 $3$ 各自本原且输入 $h$ 有定义的对，
得到 $5120$ 个局部类。它们的完整签名如表 \ref{tab:ten-signatures}，
每格表示 $(g_i,h_i')$。
条件 $5\nmid Y$ 单独检查；每个 $M_i$ 的第二行第一项为五倍数、第二项模 $5$ 非零，
共同约分又不含 $5$，因此它自动保持该条件。

\begin{table}[htbp]
\centering\small
\setlength{\tabcolsep}{4.5pt}
\caption{十元的十二种完整局部签名}\label{tab:ten-signatures}
\begin{tabular}{ccrccccccc}
\toprule
类&$h$&类数&$M_1$&$M_2$&$M_3$&$M_4$&$M_5$&$M_6$&$M_7$\\
\midrule
1&1&384&$(4,2)$&--&$(1,1)$&$(4,1)$&--&--&$(1,1)$\\
2&1&384&$(4,2)$&--&$(1,1)$&$(8,2)$&--&--&$(1,1)$\\
3&1&128&$(4,2)$&--&$(3,1)$&$(4,1)$&--&--&$(1,1)$\\
4&1&128&$(4,2)$&--&$(3,1)$&$(8,2)$&--&--&$(1,1)$\\
5&2&768&$(1,1)$&$(1,1)$&$(1,2)$&--&$(2,1)$&--&$(1,2)$\\
6&2&768&$(1,1)$&$(1,1)$&$(1,2)$&--&$(4,2)$&--&$(1,2)$\\
7&2&256&$(1,1)$&$(1,1)$&$(3,2)$&--&$(6,1)$&--&$(1,2)$\\
8&2&256&$(1,1)$&$(1,1)$&$(3,2)$&--&$(12,2)$&--&$(1,2)$\\
9&2&768&$(1,1)$&$(2,2)$&$(1,2)$&--&--&$(2,1)$&$(1,2)$\\
10&2&768&$(1,1)$&$(2,2)$&$(1,2)$&--&--&$(4,2)$&$(1,2)$\\
11&2&256&$(1,1)$&$(2,2)$&$(3,2)$&--&--&$(6,1)$&$(1,2)$\\
12&2&256&$(1,1)$&$(2,2)$&$(3,2)$&--&--&$(12,2)$&$(1,2)$\\
\bottomrule
\end{tabular}
\end{table}

\subsection{实射影覆盖}

固定一个签名。对可用矩阵 $M_i=\mat{a}{b}{c}{d}$，生成
\begin{equation}\label{eq:ten-D}
 D_i(X,Y)=h^2g_i^2(5X^2+Y^2)
 -(h_i')^2\bigl(5(aX+bY)^2+(cX+dY)^2\bigr).
\end{equation}
$D_i>0$ 等价于严格降低本原高度 $W$。
对十二组 $D_i(1,t)$ 分别执行附录 \ref{app:exact-roots} 的算法：
没有未覆盖的开区间；有限的未覆盖有理点与相应局部签名均不相容；
唯一相容的未覆盖方向为 $X=0$。
这给出引理 \ref{lem:ten-certificate} 所需的完整有限结论。
该结论可以从本附录的整数数据独立重建，无需任何发现阶段的宏搜索结果。

\section{精确验证规则与补充材料}\label{app:exact-roots}

\subsection{二次多项式正区间的有限验证}

输入是一组显式有理系数二次多项式 $p_j(t)$；本论文的输入实际为整数系数。
欲检查每个实方向是否被某个 $p_j>0$ 覆盖，使用以下精确步骤。
\begin{enumerate}
\item 删除零多项式，将每个非零多项式约成最高非零系数为正的本原形式，以识别重复根。
  符号检验仍使用原多项式，不能丢弃其原来的正负号。
\item 线性根和平方判别式的二次根用有理数表示；非平方正判别式的根，用本原不可约二次多项式及左右根标记表示。
\item 使用 Cauchy 根界及顶点分隔两个根，再以有理二分隔离不同实根。
  不同代数根距离非零，所以此过程有限终止。
\item 排序全部实根。在每个相邻根间的开区间及两端无界区间各选一个有理点。
  每个 $p_j$ 在该区间内符号不变，故一个精确符号值就决定整个区间。
\item 在每个根处，先识别哪些 $p_j$ 以它为根；其他多项式在相邻区间中的符号即为该点的符号。
  最后单独检查无穷方向，其齐次值由二次项系数给出。
\end{enumerate}
剩下的孤立有理点再代回局部签名；孤立无理方向不含非零有理向量。
七元只留下终端斜率 $0$ 和非法斜率 $2$；十元只留下终端无穷方向。

\subsection{可复现文件与数学职责}

随文源文件包包含本 \LaTeX{} 源码、编译所得 PDF，以及 \texttt{supplement/} 下的验证程序。
程序仅依赖 Python 标准库；运行时必须启用断言。
\begin{longtable}{>{\raggedright\arraybackslash}p{0.46\textwidth}>{\raggedright\arraybackslash}p{0.46\textwidth}}
\caption{补充材料的验证职责}\label{tab:supplement}\\
\toprule 文件 & 职责\\\midrule\endfirsthead
\toprule 文件 & 职责\\\midrule\endhead
\nolinkurl{verify_seven_descent_cover.py} & 重建十八个矩阵词、全部 $2916$ 个局部类、五种签名及精确实覆盖。\\
\nolinkurl{verify_eight_euclidean_descent.py} & 符号重放三个八元宏、核对坐标变换、欧几里得下降及固定终端。\\
\nolinkurl{verify_ten_descent_cover.py} & 符号重放七个十元宏、枚举全部 $5120$ 个局部类、核对十二种签名与精确实覆盖。\\
\nolinkurl{verify_ten_complete.py} & 核对四例外配对及四步入口，并调用十元完整覆盖证书。\\
\nolinkurl{verify_all_dimensions_double_triple_invariant.py} & 依正文分类选择一步；检查多素数保护、模 $2$ 仿射修复、初始化及带标签完整重放。有限测试不替代正文全称证明。\\
\nolinkurl{verify_general_invariant.py} & 独立实现的辅助闭包与路径检查；不属于七、十元有限覆盖证书的可信计算核心。\\
\nolinkurl{verify_small_dimensions.py} & 独立比较四、五元判据与完整有界整数图，核对五元级联公式、两段收缩以及成功路径的真实重放。\\
\nolinkurl{verify_paper_data.py} & 从本文源码读取矩阵、真实尺度及签名表，与精确生成的数据逐项比较。\\
\nolinkurl{run_checks.py} & 从本地补充材料统一运行上述检查；拒绝关闭断言的运行方式。\\
\bottomrule
\end{longtable}

其余几个同目录文件是精确分数、形式线性值账本或矩阵转换的共享辅助代码。
\texttt{explore\_ten\_descent.py} 在这里仅被导入三个纯算术函数；统一验证入口不运行其中的发现搜索。
七元与十元的有限覆盖检查是证明的一部分；全维数的有限数值样本仅是实现审计。
两者的证据性质不同，不应混为一谈。

\begin{thebibliography}{9}
\bibitem{mixability}
Miguel Coviello Gonzalez and Marek Chrobak.
\newblock Towards a Theory of Mixing Graphs: A Characterization of Perfect Mixability.
\newblock \emph{Theoretical Computer Science}, 845:98--121, 2020.
\newblock DOI: \href{https://doi.org/10.1016/j.tcs.2020.09.007}{10.1016/j.tcs.2020.09.007}.
\newblock 本文阅读版本：\href{https://arxiv.org/abs/1806.08875v4}{arXiv:1806.08875v4}。
\end{thebibliography}
\end{document}
